% ~9 min
\partframe{In Practice}{From Calculation to Code}

\begin{frame*}
  \frametitle{A Calculation in Rocq}
  \rocqblock{qsort2-derive.v}
  \callout<2->{14,7.6}{specification}{spec}
  \callout<3->{14,2.4}{the laws}{laws}
  \speakernote{\code{Derive}: the program starts as an unknown \\
    the spec: quicksort, then double \\
    one line of laws; Rocq fills in the program}
\end{frame*}

\begin{frame*}
  \frametitle{What Comes Out}
  \ocamlblock{qsort2.ml}
  \simplified
  \par
  \begin{onlyenv}<2->
    Hand-written (start of the talk):
    \ocamlblock{sort-times-two.ml}
  \end{onlyenv}
  \par
  \speakernote{Extracted: one pass, doubles each pivot \\
    compare: same as the hand-written one, now with a proof}
\end{frame*}

\begin{frame}
  \frametitle{Dynamic Programming: Knapsack}
  \[ K(0) = 0 \qquad K(c) = \max\,\bigl(\{0\} \cup \{\,K(c - w_i) + v_i
     \mid 0 < w_i \le c\,\}\bigr) \]
  \begin{center}
    \begin{tikzpicture}[alt={The call tree of the naive knapsack recursion
        for capacity 6 with weights 2 and 3. K(6) calls K(4) and K(3); K(4)
        calls K(2) and K(1); K(3) calls K(1) and K(0); K(2) calls K(0). The
        same capacities appear many times: K(1) twice, K(0) twice.},
      level distance = 11mm,
      level 1/.style = {sibling distance = 50mm},
      level 2/.style = {sibling distance = 24mm},
      level 3/.style = {sibling distance = 14mm},
      every node/.style = {box, font = \small}]
      \node[hot] {$K(6)$}
        child {node {$K(4)$}
          child {node {$K(2)$} child {node[draw = bad, fill = bad!8] {$K(0)$}}}
          child {node[draw = bad, fill = bad!8] {$K(1)$}}}
        child {node {$K(3)$}
          child {node[draw = bad, fill = bad!8] {$K(1)$}}
          child {node[draw = bad, fill = bad!8] {$K(0)$}}};
    \end{tikzpicture}
  \end{center}
  \uncover<2->{A hylomorphism, but it \textbf{recomputes}: exponential.}
  \speakernote{Classic DP recurrence \\
    Point at red nodes: the same capacities computed again and again \\
    naive recursion is a hylo, but exponential}
\end{frame}

\begin{frame}
  \frametitle{Remember, Don't Recompute}
  \begin{center}
    \begin{tikzpicture}[alt={A table of results for every capacity from 0
        up to c. The entry for c is computed from entries already in the
        table, for c minus w for each weight w, and is then added to the
        table.},
      every node/.style = {box, minimum width = 11mm, font = \small},
      node distance = 0mm]
      \foreach \i in {0, ..., 6} {
        \node (k\i) at (1.4 * \i, 0) {$K(\i)$};
      }
      \node[hot, minimum width = 11mm] at (k6) {$K(6)$};
      \draw[->, accent2, thick] (k4.north) to[bend left = 40]
        node[above, draw = none, fill = none, note] {read $w = 2$} (k6.north);
      \draw[->, accent2, thick] (k3.south) to[bend right = 40]
        node[below, draw = none, fill = none, note] {read $w = 3$} (k6.south);
    \end{tikzpicture}
  \end{center}
  \begin{itemize}[<+->]
    \item \textbf{Dynamorphism}: a hylomorphism that keeps a \emph{table}
      of earlier results.
    \item A law: recursion over \emph{items} $=$ recursion over
      \emph{capacities} with a table, if each step \emph{looks up} its
      subproblems.
    \item Exponential $\Longrightarrow$ one entry per capacity.
  \end{itemize}
  \slideref{R. Hinze, N. Wu, J. Gibbons. \textbf{Conjugate Hylomorphisms, or:
    The Mother of All Structured Recursion Schemes}. POPL 2015.}
  \speakernote{Keep a table of results for every capacity \\
    Point at arrows: K(6) reads K(4) and K(3) from the table \\
    dynamorphism (Hinze, Wu, Gibbons) \\
    law (\code{reenumerate}): change what we recurse on, items to capacities \\
    exponential to one entry per capacity}
\end{frame}

\begin{frame*}
  \frametitle{From Recomputation to a Table, by Rewriting}
  \begin{onlyenv}<1>
  \begin{center}
    \begin{tikzpicture}[alt={Left: the naive recursion over items, a tree
        in which K(6) calls K(4) and K(3), which call further capacities,
        with repeats. Right: the recursion over capacities, a chain from
        K(6) down to K(0); each entry looks up the entries it needs further
        down the chain instead of calling them.},
      every node/.style = {box, font = \small}]
      \begin{scope}[level distance = 10mm,
          level 1/.style = {sibling distance = 22mm},
          level 2/.style = {sibling distance = 11mm}]
        \node[hot] {$K(6)$}
          child {node {$K(4)$} child {node {$K(2)$}} child {node {$K(1)$}}}
          child {node {$K(3)$} child {node {$K(1)$}} child {node {$K(0)$}}};
        \node[draw = none, fill = none, note] at (0, 0.8)
          {over items: recompute};
      \end{scope}
      \node[draw = none, fill = none, font = \Large] at (3.4, -1)
        {$=$};
      \begin{scope}[xshift = 45mm, node distance = 3mm]
        \node[hot] (c6) at (0, 0) {$K(6)$};
        \foreach \i/\j in {5/6, 4/5, 3/4, 2/3, 1/2, 0/1} {
          \node[right = of c\j] (c\i) {$K(\i)$};
          \draw[->] (c\j) -- (c\i);
        }
        \draw[->, accent2, thick] (c6.south) to[bend right = 35]
          node[below, draw = none, fill = none, note] {look up} (c4.south);
        \draw[->, accent2, thick] (c6.north) to[bend left = 35] (c3.north);
        \node[draw = none, fill = none, note] at (3.6, 1.5)
          {over capacities: look up};
      \end{scope}
    \end{tikzpicture}
  \end{center}
  \end{onlyenv}
  \par
  \begin{onlyenv}<2->
    \rocqblock{knap-calc.v}
    \simplified
  \end{onlyenv}
  \par
  \speakernote{Left: naive recursion over items; repeats \\
    Right: count down capacities; each entry looks up earlier ones \\
    step 1: a Derive makes the algebra read through a lookup \\
    step 2: \code{reenumerate} switches the recursion; tables are persistent arrays \\
    \code{dyna\_memo}: share the table; no hand-written index arithmetic}
\end{frame*}

\begin{frame*}
  \frametitle{Knapsack: the Extracted Program}
  \ocamlblock{knap.ml}
  \simplified
  \begin{center}
    \begin{tikzpicture}[alt={The cost ladder: naive recursion is
        exponential; the calculated dynamorphism keeps one entry per
        capacity, and each lookup walks down the table; indexing the array
        directly, for constant-time lookups, is work in progress.},
      stage/.style = {box, text width = 30mm, align = center,
        minimum height = 10mm}]
      \node[stage, draw = bad, fill = bad!8] (e) {naive:\\ exponential};
      \node[stage, right = 12mm of e, draw = good, fill = good!8] (l)
        {table: one entry\\ per capacity};
      \node[stage, right = 12mm of l, dashed] (a)
        {$O(1)$ lookups:\\ in progress};
      \draw[->, thick] (e) -- node[above, note] {laws} (l);
      \draw[->, thick, dashed] (l) -- (a);
    \end{tikzpicture}
  \end{center}
  \speakernote{\code{table}: one entry per capacity, built bottom-up \\
    \code{look}: walks w steps down; the array table is persistent \\
    Ladder: exponential to one entry per capacity, all by laws \\
    Next: index the array directly for O(1) lookups (in progress)}
\end{frame*}

\begin{frame}
  \frametitle{More Calculations}
  \begin{center}
    \begin{tikzpicture}[alt={Two more examples, each from a specification
        to a calculated program. Fast exponentiation: from n
        multiplications to about log n. Maximum segment sum: from cubic to
        linear time.},
      spec/.style = {box, text width = 36mm, align = center,
        minimum height = 12mm},
      prog/.style = {box, draw = good, fill = good!8, text width = 36mm,
        align = center, minimum height = 12mm}]
      \foreach \y/\name/\s/\p/\law/\o in {
          0/{$x^n$}/{multiply $n$ times}/{$\log n$ squarings}/{deforestation}/1,
          -1.9/{Max segment sum}/{all segments: $O(n^3)$}/{one pass: $O(n)$}/{fusion, Horner, tupling}/2} {
        \uncover<\o->{
          \node[anchor = east, font = \bfseries] at (-0.3, \y) {\name};
          \node[spec, anchor = west] (s\o) at (0, \y) {\s};
          \node[prog, right = 38mm of s\o] (p\o) {\p};
          \draw[->, thick] (s\o) -- node[above, note] {\law} (p\o);
        }
      }
    \end{tikzpicture}
  \end{center}
  \speakernote{Fast power: deforest the list of binary digits \\
    max segment sum: Bird's classic, cubic to linear}
\end{frame}

\begin{frame*}
  \frametitle{Maximum Segment Sum: Spec and Result}
  \[ \op{mss} = \op{maximum} \circ \op{map}\,\op{sum} \circ \op{segments} \]
  \begin{center}
    \begin{tikzpicture}[alt={The list 3, minus 5, 2, 4, minus 1. Its
        segments include 2, 4, whose sum 6 is the largest.},
      every node/.style = {box, minimum width = 9mm}, node distance = 0mm]
      \node (a) {$3$};
      \node[right = of a] (b) {$-5$};
      \node[right = of b, hot, minimum width = 9mm] (c) {$2$};
      \node[right = of c, hot, minimum width = 9mm] (d) {$4$};
      \node[right = of d] (e) {$-1$};
      \node[draw = none, fill = none, note, right = 4mm of e]
        {best segment: $2 + 4 = 6$};
    \end{tikzpicture}
  \end{center}
  \begin{onlyenv}<2->
    \ocamlblock{mss.ml}
    \simplified
  \end{onlyenv}
  \par
  \speakernote{Spec is the definition: all segments, sum, max \\
    one pass, two accumulators; not obviously correct, but proved}
\end{frame*}

\begin{frame}
  \frametitle{Summary So Far}
  \begin{itemize}
    \item A \textbf{library}: hylomorphisms and their laws in Rocq.
    \item Optimisations are \textbf{proofs by rewriting}.
    \item Termination proved \textbf{once per split}.
    \item \textbf{No extra axioms}; OCaml \textbf{without unsafe casts}.
  \end{itemize}
  \slideref{D. Castro-Perez, M. Paviotti, M. Vollmer. \textbf{Program
    Optimisations via Hylomorphisms for Extraction of Executable Code}. ITP
    2025. Library: \url{https://github.com/dcastrop/rocq-hylomorphisms}}
  \speakernote{The library today: ITP 2025 plus later work \\
    Next: work in progress}
\end{frame}
